\section{Three problems}
\label{sec:problem}

We state three problems. One literature reaches incompatible conclusions about
a single attack. Two literatures define the same quantity and neither cites
the other. A stated threat model fails to constrain the attack it was written
for. Section~\ref{sec:problem-cause} takes the three in turn and shows that
one absence produces all of them.

\subsection{Incompatible conclusions within one literature}
\label{sec:problem-one}

Membership inference decides whether one record entered a model's training
set. Meeus et al.\ count at least \MIAnewMethods{} new methods for it proposed
in 2024 alone \cite{MSJFRM25}, and the four verdicts of
Section~\ref{sec:intro} run from success to unfalsifiability. Writing down
each paper's adversary explains the spread.
% claim: membership-inference-verdicts-disagree

Carlini et al.\ grant their adversary data drawn from the target's
distribution and the budget to train many shadow models on it. Their attack
then compares the target's per-example loss against the two resulting
distributions
\cite{CCNSTT22}. Duan et al.\ grant theirs a benchmark whose members and
non-members come from the Pile's own train and test split, which closes the
shortcut that a benchmark assembled after training leaves open
\cite{DSMMSZTCEH24}. Das et al.\ grant theirs no model access at all, and it
still wins, which places the signal in the benchmark and not in the model
\cite{DZT25}. Meeus et al.\ locate that signal precisely, reporting
\MIAwikiLo{} to \MIAwikiHi{} AUC on WikiMIA from a bag-of-words classifier whose
most predictive features are years \cite{MSJFRM25}.

Zhang et al.\ reach the strongest verdict \cite{ZDKT25}. To bound a false
positive rate, the party running the attack must know how its test statistic
is distributed under the hypothesis that the model was trained without the
record in question. Nobody can sample that hypothesis.

None of the five states a budget in a form the others could be checked
against, and the last comes closest. Zhang et al.\ give two independent
reasons that hypothesis resists sampling, and the second survives the removal
of the first. Granted full knowledge of the training set and of the training
algorithm, sampling it ``would still be prohibitively expensive as it would
require training new large models'' \cite{ZDKT25}. They declare a security
claim unfalsifiable against a party that cannot afford a training run. The
paper names that cost and moves on. It records no parameter, so nobody can
compare it with a claim made against a party of a different budget.

\subsection{The attacker-to-honest cost ratio, defined twice without contact}
\label{sec:problem-two}

A proof-of-learning asks a party claiming to have trained a model to produce a
record of the training run, which a verifier then checks \cite{JYCDTCP21}. The
scheme is secure when forging that record costs at least as much as training
honestly. Jia et al.\ state the requirement three times, as $\mathbb{E}[C_A]
\geq \mathbb{E}[C_T]$, where $C_A$ is the cost of any spoofing strategy and
$C_T$ the cost of the honest training run.
% claim: attacker-to-honest-ratio-defined-twice

The inequality carries no exponent, no adversary class, no security parameter
and no assumption. As written it holds by definition for the spoof that
retrains the model, and the authors say plainly that they cannot prove it for
anything else. Having concluded in \S VII-B3a that a passing forgery ``would
be at least as computationally expensive as what it took'' the honest trainer,
they add in the same subsection that ``given that we are not aware of a
mechanism to prove this formally, we present the following as a conjecture''
\cite{JYCDTCP21}.

Fine-grained cryptography defines the same quantity and gives it a name.
Brzuska and Couteau's Definition 6 fixes an
$(\varepsilon,\delta)$-fine-grained one-way function as one that no circuit of
size below $|f|^{1+\delta}$ inverts with probability above $\varepsilon$
\cite{BC22}. LaVigne, Lincoln and Vassilevska Williams name the exponent
itself, in their Definition 5, as ``the (relative) gap of an $T(n)$
fine-grained one-way function $f$'' \cite{LLW19}. Around that one constant
they build fine-grained key exchange and fine-grained public-key encryption.

The two literatures describe one object, and neither cites the other. The
proof-of-learning line holds the quantity and has no name for it, so it cannot
say how much more expensive forgery is. The fine-grained line holds the name
and applies it to no learning problem, so nothing in it has been instantiated
on a training run.

\subsection{Threat models stated without a resource bound}
\label{sec:problem-three}

Athalye et al.\ take the \OGpopulation{} non-certified defences accepted at
ICLR 2018 claiming white-box robustness \cite{ACW18}. Of those,
\OGobfuscated{} rely on a phenomenon the authors name obfuscated gradients.
The authors circumvent \OGbrokenFull{} completely and \OGbrokenPart{}
partially, each under the threat model its own paper stated.
% claim: threat-models-without-a-resource-bound

Two years later Tram\`er et al.\ repeat the exercise on \AAdefences{} defences
from ICLR, ICML and NeurIPS, all of which had performed adaptive evaluations
\cite{TCBM20}. They circumvent all \AAdefences{}, and they report that they
needed nothing new. In their words, ``the technical tools to evaluate defenses
properly already exist'', and the issue ``is methodological rather than
technical'' \cite{TCBM20}.

A defence and its break can both be correct. The threat models in question fix
the perturbation set and the access level, and they leave open what the
attacker may compute and what counts as success. A model that fixes neither of
those constrains no attack, so a defence can satisfy it while an attack
violates the defence, with neither party in error.

\subsection{The common cause}
\label{sec:problem-cause}

One cause produces all three problems. Most of these works state no budget for
any party, so nobody has a position to hold. Researchers who disagreed about
what an adversary is would already be an improvement, because two positions
can be compared. Section~\ref{sec:findings} reports how many.

Take the three in turn. In the first, the five membership inference results
grant their adversaries different things, and each verdict follows its grant.
Carlini et al.\ grant the target's distribution and the compute to train
shadow models on it. Duan et al.\ grant a benchmark that closes the shortcut.
Das et al.\ grant no model access at all. Nobody can order the results,
because no two of them price the same party in the same unit.

In the second, the proof-of-learning line holds the quantity and has no name
for it, while the fine-grained line holds the name and has instantiated it on
no training run. Neither finds the other by searching, because they share no
vocabulary. The object both describe is a relation between two parties' work,
and only one of the two literatures writes it as an exponent.

In the third, the threat models fix the perturbation set and the level of
access, and they leave open what the attacker may compute and what counts as
success. A defence can satisfy such a model and an attack can then break the
defence, with neither party in error, because the model never ruled the attack
out.

The field has a name for the third case and for neither of the first two. Arp
et al.'s tenth pitfall, quoted in Section~\ref{sec:intro}, is the threat
model, defined as the case where ``the security of machine learning is not
considered, exposing the system to a variety of attacks, such as poisoning and
evasion attacks'' \cite{AQPWPWCR22}. That pitfall and its recommendations
quantify the adversary by no computation, memory, sample or query count. The
field diagnoses its own threat models as unstated or imprecise, and never as
unpriced.
